\documentclass[12pt]{article}
\usepackage{amsmath,amssymb}

\begin{document}
\section*{{2020 AMC 10B Problems/Problem 25}}
Source: AoPS Wiki

\subsection*{Solution}
Note that $96 = 2^5 \cdot 3$ . Since there are at most six not necessarily distinct factors $>1$ multiplying to $96$ , we have six cases: $k=1, 2, ..., 6.$ Now we look at each of the six cases. $k=1$ : We see that there is $1$ way, merely $96$ . $k=2$ : This way, we have the $3$ in one slot and $2$ in another, and symmetry. The four other $2$ 's leave us with $5$ ways and symmetry doubles us so we have $10$ . $k=3$ : We have $3, 2, 2$ as our baseline. We need to multiply by $2$ in $3$ places, and see that we can split the remaining three powers of $2$ in a manner that is $3-0-0$ , $2-1-0$ or $1-1-1$ . A $3-0-0$ split has $6 + 3 = 9$ ways of happening ( $24-2-2$ and symmetry; $2-3-16$ and symmetry), a $2-1-0$ split has $6 \cdot 3 = 18$ ways of happening (due to all being distinct) and a $1-1-1$ split has $3$ ways of happening ( $6-4-4$ and symmetry) so in this case we have $9+18+3=30$ ways. $k=4$ : We have $3, 2, 2, 2$ as our baseline, and for the two other $2$ 's, we have a $2-0-0-0$ or $1-1-0-0$ split. The former grants us $4+12=16$ ways ( $12-2-2-2$ and symmetry and $3-8-2-2$ and symmetry) and the latter grants us also $12+12=24$ ways ( $6-4-2-2$ and symmetry and $3-4-4-2$ and symmetry) for a total of $16+24=40$ ways. $k=5$ : We have $3, 2, 2, 2, 2$ as our baseline and one place to put the last two: on another two or on the three. On the three gives us $5$ ways due to symmetry and on another two gives us $5 \cdot 4 = 20$ ways due to symmetry. Thus, we have $5+20=25$ ways. $k=6$ : We have $3, 2, 2, 2, 2, 2$ and symmetry and no more twos to multiply, so by symmetry, we have $6$ ways. Thus, adding, we have $1+10+30+40+25+6=\boxed{\textbf{(A) } 112}$ . ~kevinmathz

\end{document}
